Para una exposición universitaria de **20 minutos** sobre **qubits fotónicos**, asumiendo que tu audiencia ya conoce los fundamentos de la computación cuántica (postulados, operaciones tensoriales, compuertas y algoritmos), te recomiendo una estructura **concisa, técnica y enfocada en lo distintivo de los qubits fotónicos**. Aquí tienes una posible división:

---

### **1. Introducción (2 min)**
- **Objetivo**: Contextualizar por qué los qubits fotónicos son relevantes.  
  - Frase inicial: *"Si los qubits tradicionales (superconductores, iones) son como autos, los fotónicos son como aviones: no compiten, resuelven problemas distintos".*  
  - Ventajas clave: Baja decoherencia, transmisión a distancia, operación a temperatura ambiente.  
  - Desafío principal: *"¿Cómo hacer computación cuántica sin interacción fuerte entre fotones?"*.

---

### **2. Codificación de información en qubits fotónicos (4 min)**
- **Enfatiza las diferencias** con qubits "ordinarios":  
  - **Degrees of freedom**: Polarización (|H⟩/|V⟩), momento angular orbital (OAM), tiempo-bin, número de fotones.  
  - **Ejemplo concreto**: Mostrar cómo un CNOT se implementa en polarización usando cristales no lineales (BBO) o interferómetros.  
  - **Gráfico útil**: Tabla comparativa de encodings y sus ventajas (ej: polarización para QKD, OAM para alta dimensión).  

---

### **3. Operaciones y Compuertas (4 min)**
- **Problema central**: Fotones casi no interactúan → Soluciones ingeniosas.  
  - **Medición indirecta (MBQC)**: Ejemplo con cluster states (Universidad de Bristol).  
  - **No linealidades ópticas**: Kerr effect en anillos de silicio (investigación reciente, ej: Nature Photonics 2023).  
  - **Puertas deterministas vs probabilistas**: Ventaja de los esquemas basados en medición.  
  - **Demo visual**: Circuito óptico de una compuerta CZ usando acopladores direccionales (ideal si puedes mostrar un chip fotónico integrado).  

---

### **4. Aplicaciones Prácticas (4 min)**
- **Enfócate en lo único**:  
  - **Ventaja en comunicación**: QKD (BB84 con polarización), teleportación (experimento de Zeilinger).  
  - **Ventaja en algoritmos**: Boson Sampling (¿por qué es difícil simularlo clásicamente?).  
  - **Híbridos**: Fotones como "mensajeros" entre qubits superconductores/átomos (ej: Quantum Internet).  

---

### **5. Estado Actual y Retos (3 min)**
- **Escalabilidad**:  
  - Límites de pérdidas en fibra vs. integración fotónica (PsiQuantum, Xanadu).  
  - Tasas de error y corrección (códigos ópticos como el "cat code").  
- **Hardware**:  
  - Fuentes de fotones únicos (quantum dots vs SPDC).  
  - Detectores superconductores (SNSPD) vs. SPAD.  

---

### **6. Cierre Inspirador (2 min)**
- **Futuro**:  
  - *"La fotónica no compite con otras plataformas: es la llave para la red cuántica y algoritmos específicos"*.  
  - **Paper clave para profundizar**: *"Photonic quantum information processing"* (Walmsley, Science 2019).  
- **Pregunta retórica**: *"Si los fotones son tan buenos para transmitir información, ¿serán la columna vertebral de la Internet cuántica?"*.  

---

### **Tips Adicionales**  
- **Diapositivas**: Máximo 10-12 slides (1 por minuto, pero algunas son solo imágenes/gráficos).  
- **Ejemplo real**: Menciona un experimento reciente (ej: teleportación de 3 estados dimensionales con OAM, 2023).  
- **Interacción**: Lanza una pregunta al público: *"¿Creen que la corrección de errores en fotónica será suficiente para competir con superconductores?"*.  

Si quieres profundizar en algún punto (ej: detalles de Boson Sampling o cómo funciona un interferómetro de Hong-Ou-Mandel), dime y ajustamos. ¡Éxito en tu exposición! ��